% Discussion

\section{Discussion}
\label{sec:discussion}

\paragraph{Execution flywheel: every run refines the shared cache.}
A key property of \sysname is that the shared Function cache improves with natural usage---no dedicated exploration phase is required.
Each accepted run is written back through the same write-back path: \textsc{Omni} projects the new observation into the existing UDEG node family, adds the run as additional evidence, and refines the Function's precondition and action evidence offline.
The consequence is that the same UDEG node accumulates observations from diverse devices, users, app versions, and screen configurations over time.
A Function that was originally recorded on one device and screen size gradually acquires anchor evidence from foldable displays, different OEM skins, and varying font scales.
This makes subsequent cache hits more robust---not because the system was trained on more data, but because the PageVector and anchor evidence naturally span a wider layout manifold.
The cost structure follows: the per-Function online model cost is paid once, on the first execution that discovers the Function.
Every subsequent cache hit amortizes that cost across all future users.
In workloads where the same recurring Function is triggered by many users, the marginal online model cost per invocation approaches zero as the cache matures.
This flywheel is absent in existing memory-based systems that build per-user, per-app memory from dedicated exploration: each new user resets the cold-start cost, and the memory does not cross-fertilize across users or devices.

Our efficient page matching and anchor relocation algorithms make repeated write-back practical at runtime. As a result, every successful execution can be incorporated into the shared cache with little additional overhead, allowing execution knowledge to accumulate continuously through normal usage.